\documentclass[12pt]{article}
\usepackage[utf8]{inputenc}
\usepackage[T1]{fontenc}
\usepackage{hyperref}
\title{Ponosakan Language Resources for K-12 Curriculum Development}
\author{Compiled by Claude Sonnet 5.5 (claude-sonnet-5-5) for Andr\'as Kornai}
\date{}
\begin{document}
\maketitle
\section{Introduction}
Ponosakan (ISO 639-3 pns) is a moribund Austronesian language near Belang, Southeast Minahasa, North Sulawesi, Indonesia. It is often called Minahasan locally, but scholars place it in the Gorontalo--Mongondow group. Wikipedia reports three fluent speakers as of November 2024 and says local revival efforts, including teaching in primary schools, began in 2024. This report was not verified beyond the search snippets.
\section{Primary Educational Resources}
No specific textbooks were found. Wikipedia mentions school teaching started in 2024, but the materials were not identified.
\section{Linguistic and Reference Materials}
\begin{itemize}
\item Wikipedia, Ponosakan language. \url{https://en.wikipedia.org/wiki/Ponosakan_language}
\item Omniglot, Ponosakan (overview and history of decline). \url{https://omniglot.com/writing/ponosakan.htm}
\item An article from the Language Documentation and Description journal that appeared in search results; its relevance to Ponosakan was not confirmed. \url{https://lddjournal.org/article/1203/galley/2447/download}
\end{itemize}
\section{Religious and Community Resources}
\begin{itemize}
\item Joshua Project, Ponasakan in Indonesia (profile). \url{https://joshuaproject.net/people_groups/14440}
\end{itemize}
\section{Digital and Audio Resources}
None found.
\section{Curriculum Development Context}
With very few elderly speakers, a heritage curriculum would depend on urgent recording of the remaining speakers. Documentation found online is minimal, and contact with the local education authorities in Southeast Minahasa would be needed to learn what the 2024 school teaching consists of.
\end{document}